\chapter{Las ciencias de la vida}
\label{ch:life}

\section{El único sector que ya posee algo}

En todo el resto de la Parte III el problema es construir una capacidad. Aquí la capacidad
existe, tiene cuarenta años, tiene productos, tiene reputación, y nunca se ha organizado
para vender nada.

La biotecnología cubana empezó a comienzos de los años ochenta como una decisión estatal
deliberada de construir capacidad de investigación en un país pobre, y produjo, contra toda
expectativa razonable, una capacidad genuina. Las instituciones ---el Centro de Ingeniería
Genética y Biotecnología, el Instituto Finlay, el Centro de Inmunología Molecular y las
demás, consolidadas desde 2012 en el holding BioCubaFarma--- han entregado productos que no
eran derivativos.

La vacuna contra el meningococo B desarrollada en el Instituto Finlay en los años ochenta
fue la primera de su tipo en cualquier parte. La vacuna contra el \emph{Haemophilus
influenzae} tipo b licenciada en 2004 usó un antígeno enteramente sintético, la primera
vacuna de antígeno sintético en uso humano rutinario. El Heberprot-P, un factor de
crecimiento epidérmico recombinante que se inyecta en las úlceras del pie diabético,
atiende una condición que causa una parte grande de las amputaciones del mundo y no tiene
equivalente cercano en otra parte. El CIMAvax-EGF, una vacuna terapéutica contra el cáncer
de pulmón, ha resultado lo bastante interesante para los oncólogos norteamericanos como
para que un centro oncológico de Estados Unidos persiguiera ensayos con ella al amparo de
las excepciones de licencia. Y en 2021 el sector desarrolló, ensayó, fabricó y desplegó sus
propias vacunas contra la COVID-19, vacunando a más del noventa por ciento de la población
---único país de América Latina que lo hizo, y lo hizo sin poder importar jeringuillas de
forma fiable---.\unv{}

¿Por qué funcionó esto cuando no funcionó nada más? Tres razones, y vale la pena nombrarlas
porque la reforma no debe destruirlas.

Era un \emph{circuito cerrado}. Investigación, desarrollo clínico, fabricación y despliegue
estaban dentro de un solo sistema con un solo cliente, el sistema nacional de salud, de modo
que un producto pasaba del laboratorio a la población sin un mercado intermedio y sin una
negociación de reembolso. Es una ventaja estructural genuina para desarrollar productos
dirigidos a una población entera.

Tenía \emph{capital paciente}. El Estado lo financió durante una década antes de que hubiera
resultados, cosa que ningún accionista habría hecho, y no lo cerró durante el Período
Especial.

Y no tenía \emph{restricción de rentabilidad en la elección de objetivos}. Trabajó sobre la
meningitis B, sobre el pie diabético, sobre el cáncer de pulmón en fumadores ---condiciones
de interés para Cuba y para los países pobres en general antes que para los mercados que
fijan las prioridades farmacéuticas globales---.

\section{Por qué no gana casi nada}

Y sin embargo. Las exportaciones farmacéuticas y biotecnológicas cubanas son, por cualquier
medida publicada, una fracción pequeña de lo que debería valer una cartera así ---del orden
de unos pocos cientos de millones de dólares al año, frente a un mercado global de billones
y frente a unas exportaciones cubanas de \emph{servicios} médicos que en su punto máximo
fueron veinte o treinta veces mayores---.\unv{}

Las razones no son científicas. Son regulatorias, contractuales y financieras, y cada una
tiene arreglo.

\emph{Ninguna conversión regulatoria.} Casi nada de la cartera tiene aprobación de la
Administración de Alimentos y Medicamentos de Estados Unidos ni de la Agencia Europea de
Medicamentos, y buena parte de la evidencia clínica no se generó según los estándares que
esas agencias exigen ---no porque los ensayos fueran deshonestos, sino porque se diseñaron
para el registro cubano y se publicaron en revistas cubanas---. Un producto sin aprobación de
una autoridad rigurosa no puede venderse en los mercados que pagan, no puede licenciarse a un
socio que venda allí, y no puede cobrarse a nada parecido a su valor terapéutico. Las vacunas
contra la COVID son el caso más claro: Cuba desarrolló vacunas que funcionaban y no pudo
venderlas, porque no completaron la vía internacional de precalificación.

\emph{Ningún socio de distribución.} El valor en la industria farmacéutica se captura a
través de los asuntos regulatorios, la distribución, la fuerza de ventas y la negociación de
reembolso ---capacidades que Cuba no tiene y que tardaría veinte años en construir---. La
respuesta es asociarse, y Cuba en gran medida no se ha asociado, en parte por las dos razones
que siguen.

\emph{Ningún capital.} Llevar un solo producto a través de ensayos internacionales cuesta más
que el presupuesto anual de todo el sector, y un soberano en impago no puede financiarlo.

\emph{Una estructura de propiedad con la que nadie quiere contratar.} Una empresa
farmacéutica extranjera que entra en un acuerdo de licencia está comprando un flujo de
obligaciones exigibles durante quince años. Su contraparte en Cuba es un holding estatal, en
una jurisdicción sin tribunales independientes, sujeta a un estatuto que expone a quien
trafique con propiedad confiscada a litigio en Estados Unidos, con un soberano que congeló
las cuentas de empresas extranjeras en 2009. La ciencia no es el riesgo.

\emph{Y el regulador es el productor.} El regulador cubano de medicamentos, el CECMED, es un
organismo estatal que regula productos fabricados por organismos estatales, dentro de un
único sistema que responde a un solo gobierno. Cualquiera que sea la integridad de su
personal ---y su reputación técnica no es mala--- estructuralmente no es independiente, y la
independencia estructural es precisamente lo que busca un regulador extranjero o un socio
extranjero. Es el mismo conflicto de interés que el capítulo~\ref{ch:energy} identificó en la
electricidad, donde el operador, el propietario, el planificador y el regulador son una sola
entidad.

\section{El diseño}

La estrategia es convertir un activo científico existente en un sector regulado, invertible y
exportador sin destruir el circuito cerrado que lo hizo funcionar. El orden importa: primero
el regulador, porque nada más es posible hasta que una contraparte extranjera pueda confiar en
algo.

\begin{design}{Un regulador de medicamentos independiente}
\label{d:cecmed}
\begin{rules}
\rl{1}{Separación} El regulador de medicamentos se constituye como autoridad independiente
con las garantías de mandato, presupuesto y publicación del Diseño~\ref{d:courts},
institucionalmente separada de BioCubaFarma, del ministerio de salud en su calidad de
propietario, y de toda entidad productora.

\rl{2}{Financiado por tasas y publicado} Financiado con tasas de solicitud a aranceles
publicados, con todos los informes de evaluación, aprobaciones, denegaciones y hallazgos de
inspección publicados, como publican los suyos las agencias europea y norteamericana. La
publicación es lo que hace verificables desde fuera los juicios de un regulador, que es el
activo entero que se está construyendo.

\rl{3}{Meta: condición de autoridad rigurosa} Un programa explícito y con fecha para alcanzar
la condición de autoridad reconocida por la Organización Mundial de la Salud, con los
criterios de evaluación, el análisis de brechas y el avance publicados anualmente. Es un
proyecto técnico de varios años con un punto de llegada definido, y es la inversión
institucional de mayor rendimiento disponible para Cuba en cualquier sector.

\rl{4}{Reconocimiento entretanto} Reconocimiento automático de las aprobaciones de reguladores
extranjeros rigurosos para los productos que Cuba importa, de modo que la escasa capacidad de
evaluación se dedique a los productos cubanos que buscan exportar y no a revisar de nuevo
medicamentos ya aprobados en otras partes. Esto además acorta de inmediato la cola de
medicinas importadas que hoy faltan en las farmacias.

\rl{5}{Ensayos con estándar internacional} Todos los ensayos clínicos cubanos conducidos y
registrados conforme a estándares internacionales de buenas prácticas clínicas, en registros
internacionales, con monitoreo y revisión ética independientes. Los datos de ensayos cubanos
se descuentan hoy en el exterior por razones de procedimiento antes que de sustancia, y esta
cláusula es lo que vuelve vendibles cuarenta años de ciencia acumulada.
\end{rules}
\end{design}

\begin{design}{Capital y asociación}
\label{d:biopartner}
\begin{rules}
\rl{1}{Empresas conjuntas a nivel de producto} Asociación al nivel del producto individual
---un activo, un socio, territorios definidos, hitos definidos--- antes que al nivel de la
institución. Esto levanta capital y compra capacidad regulatoria y de distribución sin vender
la base investigadora.

\rl{2}{Conservar la plataforma, licenciar el producto} Cuba conserva la propiedad de sus
instituciones de investigación, su fabricación y sus plataformas tecnológicas subyacentes, y
licencia derechos sobre productos específicos en mercados específicos. Los comparadores son
Singapur e Irlanda ---países que construyeron sectores farmacéuticos alojando y asociándose
antes que originando--- y no Suiza, que no es un modelo disponible para un país de nueve
millones sin capital.

\rl{3}{Términos exigibles} Contratos bajo derecho extranjero con arbitraje internacional,
según el capítulo~\ref{ch:capital}. Un acuerdo de licencia justiciable solo ante tribunales
cubanos vale una fracción del mismo acuerdo justiciable en Londres o Singapur, y el descuento
es mayor que cualquier cosa que Cuba pudiera negociar en la regalía.

\rl{4}{Científicos de la diáspora} Los científicos cubanos en el exterior son el puente
natural hacia socios, reguladores y capital extranjeros, y el Diseño~\ref{d:diaspora}.5 se
aplica directamente. Varias de las personas mejor situadas para llevar un producto cubano a
través de una solicitud ante la FDA son cubanos que se fueron.

\rl{5}{El abasto doméstico primero} Nada de este diseño permite exportar un medicamento que no
esté disponible en las farmacias cubanas. El abasto doméstico es la primera prioridad de la
producción, y se enuncia como regla porque la presión de los ingresos irá en la dirección
contraria.
\end{rules}
\end{design}

\section{Las misiones médicas, reformadas}

La mayor exportación de Cuba durante más de una década fue el trabajo de sus profesionales de
la salud, y el capítulo~\ref{ch:venezuela} informó de sus dos mitades: un servicio real
realmente prestado en lugares que ningún otro país cubría, financiado pagando a los
trabajadores una fracción del valor del contrato y, en varios destinos, quedándose con sus
pasaportes.

El programa debería conservarse. Es valioso, es genuinamente bueno para los países
receptores, y es una de las pocas cosas que Cuba vende que no vende nadie más. Pero el arreglo
laboral tiene que cambiar, y el argumento para cambiarlo no es solo ético. Es que la
restricción que hoy limita a la industria es la oferta de médicos dispuestos ---los
profesionales de la salud cubanos se están yendo del país por completo, y un programa que los
trata como reclutas los pierde hacia la emigración, que no le rinde nada a Cuba---.

\begin{design}{La exportación de servicios médicos}
\label{d:missions}
\begin{rules}
\rl{1}{Contratos individuales} Cada trabajador de la salud contrata individualmente, con el
valor completo del contrato revelado y una proporción mínima publicada pagadera al trabajador.
El Estado puede retener una parte ---hay una pretensión pública legítima sobre el rendimiento
de una formación financiada públicamente--- y esa parte se publica, se legisla, y es la misma
para todos.

\rl{2}{Ningún control de documentos ni de movimiento} Los pasaportes los guardan sus dueños.
Se suprimen las restricciones de residencia, asociación y relaciones personales. La sanción
por dejar una misión es contractual y no una prohibición de regresar al propio país, y se
aplica el Diseño~\ref{d:diaspora}.2.

\rl{3}{Voluntario} La participación es voluntaria, la negativa no acarrea consecuencia
profesional, y esto es auditable.

\rl{4}{Conforme a estándares laborales} El programa se pone en conformidad con los convenios
internacionales del trabajo, y la conformidad se verifica externamente y se publica. Es lo que
convierte un programa bajo crítica internacional permanente en un servicio exportable,
ampliable y defendible.

\rl{5}{Piso de dotación doméstica} Un nivel mínimo publicado de dotación del sistema de salud
cubano, por debajo del cual no se dotan misiones. Según el capítulo~\ref{ch:premises}, la
gente es la escasez vinculante, y el sistema doméstico es el activo que el libro entero está
escrito para proteger.
\end{rules}
\end{design}

\section{Dos líneas de negocio más}

\emph{Los ensayos clínicos como exportación.} Cuba tiene atributos por los que pagan las
organizaciones de investigación por contrato: una población con cobertura sanitaria universal
y expedientes médicos completos, alto seguimiento y retención porque la gente no se muda ni
pierde su seguro, una red de atención primaria densa en médicos y organizada por población
definida, y una base de costos de ensayo muy por debajo de Norteamérica o Europa. Bajo el
Diseño~\ref{d:cecmed}.5 y un marco ético independiente, correr fases de ensayos
internacionales en Cuba es una exportación genuina, poco intensiva en capital y de alta
calificación, que además le da al sistema doméstico acceso a tratamientos en investigación.
Exige estándares escrupulosos de consentimiento ---una población con acceso limitado a
medicinas es una población a la que se puede inducir a enrolarse--- y la revisión ética tiene
que ser independiente tanto del patrocinador como del Estado, que es una exigencia más fuerte
que la que aplica la mayoría de los países de ingreso medio que albergan ensayos.

\emph{El turismo de salud.} Cuba lleva décadas tratando pacientes extranjeros y podría hacer
mucho más, en cardiología, ortopedia, oftalmología, oncología, rehabilitación neurológica y
tratamiento de adicciones, a precios muy por debajo de Estados Unidos y con resultados que
resisten la comparación. La restricción es de nuevo la del capítulo~\ref{ch:premises}: cada
paciente extranjero consume tiempo de un clínico cubano. El Diseño~\ref{d:delivery}.3 ya
obliga: instalaciones designadas, personal acotado, ingresos al sistema público, y el programa
se detiene si el acotamiento no puede sostenerse. Este capítulo no añade nada a eso salvo la
observación de que el turismo de salud es el visitante de mayor valor por noche de cama
disponible bajo la métrica del capítulo~\ref{ch:tourism}, y el menos dañino ambientalmente.

\section{Lo que se está protegiendo}

Una advertencia final, porque las reformas de arriba van todas en dirección de los mercados y
la asociación, y esa dirección tiene un modo de fallo conocido.

El sector biotecnológico cubano funciona porque es un circuito cerrado con un cliente público,
capital paciente y ninguna obligación de elegir objetivos por tamaño de mercado. Ábralo mal a
la asociación y se convierte en un fabricante por contrato de la cartera de otro, con su
agenda de investigación fijada en Nueva Jersey y su función de salud pública abandonada
---momento en el cual Cuba habrá vendido su única capacidad originadora a cambio de un margen
de fabricación---.

Las cláusulas que protegen contra eso son el Diseño~\ref{d:biopartner}.2, que conserva la
plataforma y licencia solo productos; el Diseño~\ref{d:biopartner}.5, que pone la farmacia
doméstica primero; y el presupuesto público de investigación, que protege el piso del
capítulo~\ref{ch:welfare}. Son toda la protección que hay, no se aplican solas, y a un
gobierno futuro bajo presión fiscal le ofrecerán mucho dinero por debilitarlas.
